\section{What Can Be Identified?}
Accurate prediction of detection does not by itself establish accurate
recovery of fraud onset. The following restricted example already exhibits
ambiguity.

\begin{proposition}
Suppose $T_f\sim\mathrm{Exp}(\alpha)$ and the independent detection delay
$L=T_d-T_f\sim\mathrm{Exp}(\beta)$, with positive unequal rates. The
distribution of $T_d$ alone cannot distinguish $(\alpha,\beta)$ from
$(\beta,\alpha)$.
\end{proposition}
\begin{proof}
The convolution density is
\begin{equation}
 p_d(t)=\frac{\alpha\beta}{\beta-\alpha}
       \left(e^{-\alpha t}-e^{-\beta t}\right).
\end{equation}
Exchanging the rates changes the signs of both the denominator and the
parenthesized difference, leaving the density unchanged. Detection survival
is also unchanged, so administrative right-censoring cannot resolve the
ambiguity. The mean fraud onset, $1/\alpha$, generally changes.
\end{proof}

This is a counterexample, not a general identification theorem. It shows
that detection-time fit alone is insufficient even in a simple submodel.
Flexible neural hazards can introduce further ambiguity. Restricting a
hazard or regularizing parameters selects solutions but does not establish
that the selected onset corresponds to the true process.

Possible additional information includes an audited subset with onset labels
or intervals, externally estimated detection delays, or covariates with
justified effects on only one process. Each route requires explicit
assumptions and sensitivity analysis. We will first compare unconstrained
estimation with known-delay and partially audited settings. A meaningful
outcome may be calibrated onset intervals or sensitivity bounds rather than
point identification.
